
\documentclass{article}
\usepackage{colortbl}
\usepackage{makecell}
\usepackage{multirow}
\usepackage{supertabular}

\begin{document}

\newcounter{utterance}

\centering \large Interaction Transcript for game `clean\_up', experiment `2\_hard\_3obj\_ru', episode 3 with deepseek{-}v4{-}pro.
\vspace{24pt}

{ \footnotesize  \setcounter{utterance}{1}
\setlength{\tabcolsep}{0pt}
\begin{supertabular}{c@{$\;$}|p{.15\linewidth}@{}p{.15\linewidth}p{.15\linewidth}p{.15\linewidth}p{.15\linewidth}p{.15\linewidth}}
    \# & $\;$A & \multicolumn{4}{c}{Game Master} & $\;\:$B\\
    \hline

    \theutterance \stepcounter{utterance}  
    & & \multicolumn{4}{p{0.6\linewidth}}{
        \cellcolor[rgb]{0.9,0.9,0.9}{
            \makecell[{{p{\linewidth}}}]{
                \texttt{\tiny{[P1$\langle$GM]}}
                \texttt{Я твой ведущий, и ты играешь в совместную игру, где игровым полем служит следующая сетка:} \\
\\ 
\texttt{```} \\
\texttt{ 1234567} \\
\texttt{╔═╤═╤═══╗} \\
\texttt{║L│◌│◌◌◌║ 1} \\
\texttt{║C├─┤◌◌◌║ 2} \\
\texttt{╟─┤◌├─┬─╢ 3} \\
\texttt{║◌│◌│◌│◌║ 4} \\
\texttt{╟─┼─┘◌│◌║ 5} \\
\texttt{║◌│◌◌┌┴─╢ 6} \\
\texttt{║◌│◌◌│◌P║ 7} \\
\texttt{╚═╧══╧══╝} \\
\\ 
\texttt{```} \\
\\ 
\texttt{* По верхнему краю показаны x{-}координаты, увеличивающиеся вправо, а по правому краю — y{-}координаты, увеличивающиеся вниз.} \\
\texttt{* На твоей сетке случайным образом размещены следующие объекты: 'C', 'L', 'P'.} \\
\\ 
\\ 
\texttt{Другой игрок видит вариант игрового поля, где объекты размещены в других случайных местах. Ты не видишь поле другого игрока, и он не видит твоё.} \\
\\ 
\\ 
\texttt{**Цель:**} \\
\\ 
\texttt{Оба игрока должны перемещать объекты на своём фоне так, чтобы одинаковые объекты оказались в одних и тех же координатах. Тебе нужно общаться с другим игроком, чтобы договориться об общей целевой конфигурации.} \\
\\ 
\\ 
\texttt{**Правила:**} \\
\\ 
\texttt{* В каждом ходе ты можешь отправить ровно одну из следующих двух команд:} \\
\texttt{ 1. `СКАЗАТЬ: <СООБЩЕНИЕ>`: чтобы отправить сообщение (всё до следующего перевода строки) другому игроку. Я перешлю его твоему партнёру.} \\
\texttt{2. `ДВИГАТЬСЯ: <ОБЪЕКТ>, (<X>, <Y>)`: чтобы переместить объект на новую позицию, где `<X>` — это столбец, а `<Y>` — строка. Я сообщу тебе, был ли ход успешным или нет.} \\
\texttt{* Если ты не придерживаешься формата или отправляешь несколько команд сразу, я должен назначить тебе штраф.} \\
\texttt{* Если оба игрока наберут больше чем 8 штрафов, вы оба проиграете.} \\
\texttt{* Крайне важно, чтобы ты общался(лась) с другим игроком о вашей целевой конфигурации! *Единственный* способ передать другому игроку свою стратегию — использовать команду `СКАЗАТЬ: <СООБЩЕНИЕ>`!} \\
\\ 
\\ 
\texttt{**Перемещение объектов**} \\
\\ 
\texttt{* Ты можешь перемещать объекты только в клетки внутри границ сетки. Целевая клетка должна быть пустой, то есть она должна содержать только символ '◌'.} \\
\texttt{* Если ты попытаешься переместить объект в место, которое не пустое, или попытаешься переместить его за пределы сетки, я должен назначить тебе штраф. У тебя будет ещё одна попытка.} \\
\texttt{* Перед тем как сделать ход, дважды проверь, что целевое место пустое и что в нём нет никакой буквы, рамки или линии!} \\
\\ 
\\ 
\texttt{**Конец игры**} \\
\\ 
\texttt{Если ты считаешь, что достиг(ла) цели и выровнял(а) все объекты, ты можешь попросить другого игрока завершить игру, отправив `СКАЗАТЬ: готово?`. Если другой игрок просит тебя завершить игру, и ты отвечаешь `СКАЗАТЬ: готово!`, игра завершится.} \\
\\ 
\texttt{Оба игрока выигрывают, если игра закончится в пределах 12 раундов, где один раунд определяется как то, что оба игрока отправили по одной корректной команде.} \\
\\ 
\texttt{**Подсчёт очков:**} \\
\\ 
\texttt{Чем ближе одинаковые объекты расположены на обоих игровых полях, тем больше очков ты получаешь. Штрафы уменьшают твои очки. Сможешь побить рекорд?} \\
\\ 
\\ 
\texttt{Пожалуйста, отправь сообщение другому игроку, чтобы начать игру!} \\
            }
        }
    }
    & \\ \\

    \theutterance \stepcounter{utterance}  
    & \multicolumn{4}{p{0.6\linewidth}}{
        \cellcolor[rgb]{0.9,0.9,0.9}{
            \makecell[{{p{\linewidth}}}]{
                \texttt{\tiny{[P1$\rangle$GM]}}
                \texttt{СКАЗАТЬ: Привет! Давай договоримся о целевой конфигурации. Какие у тебя координаты объектов C, L, P на твоём поле? У меня: C в (1,2), L в (1,1), P в (7,7).} \\
            }
        }
    }
    & & \\ \\

    \theutterance \stepcounter{utterance}  
    & & & \multicolumn{2}{p{0.3\linewidth}}{
        \cellcolor[rgb]{0.9,0.9,0.9}{
            \makecell[{{p{\linewidth}}}]{
                \texttt{\tiny{[GM$|$GM]}}
                \texttt{СКАЗАТЬ: Привет! Давай договоримся о целевой конфигурации. Какие у тебя координаты объектов C, L, P на твоём поле? У меня: C в (1,2), L в (1,1), P в (7,7).} \\
            }
        }
    }
    & & \\ \\

    \theutterance \stepcounter{utterance}  
    & & \multicolumn{4}{p{0.6\linewidth}}{
        \cellcolor[rgb]{0.9,0.9,0.9}{
            \makecell[{{p{\linewidth}}}]{
                \texttt{\tiny{[GM$\rangle$P2]}}
                \texttt{Я твой ведущий, и ты играешь в совместную игру, где игровым полем служит следующая сетка:} \\
\\ 
\texttt{```} \\
\texttt{ 1234567} \\
\texttt{╔═╤═╤═══╗} \\
\texttt{║◌│◌│◌◌◌║ 1} \\
\texttt{║◌├─┤◌L◌║ 2} \\
\texttt{╟─┤P├─┬─╢ 3} \\
\texttt{║◌│◌│◌│◌║ 4} \\
\texttt{╟─┼─┘C│◌║ 5} \\
\texttt{║◌│◌◌┌┴─╢ 6} \\
\texttt{║◌│◌◌│◌◌║ 7} \\
\texttt{╚═╧══╧══╝} \\
\\ 
\texttt{```} \\
\\ 
\texttt{* По верхнему краю показаны x{-}координаты, увеличивающиеся вправо, а по правому краю — y{-}координаты, увеличивающиеся вниз.} \\
\texttt{* На твоей сетке случайным образом размещены следующие объекты: 'C', 'L', 'P'.} \\
\\ 
\\ 
\texttt{Другой игрок видит вариант игрового поля, где объекты размещены в других случайных местах. Ты не видишь поле другого игрока, и он не видит твоё.} \\
\\ 
\\ 
\texttt{**Цель:**} \\
\\ 
\texttt{Оба игрока должны перемещать объекты на своём фоне так, чтобы одинаковые объекты оказались в одних и тех же координатах. Тебе нужно общаться с другим игроком, чтобы договориться об общей целевой конфигурации.} \\
\\ 
\\ 
\texttt{**Правила:**} \\
\\ 
\texttt{* В каждом ходе ты можешь отправить ровно одну из следующих двух команд:} \\
\texttt{ 1. `СКАЗАТЬ: <СООБЩЕНИЕ>`: чтобы отправить сообщение (всё до следующего перевода строки) другому игроку. Я перешлю его твоему партнёру.} \\
\texttt{2. `ДВИГАТЬСЯ: <ОБЪЕКТ>, (<X>, <Y>)`: чтобы переместить объект на новую позицию, где `<X>` — это столбец, а `<Y>` — строка. Я сообщу тебе, был ли ход успешным или нет.} \\
\texttt{* Если ты не придерживаешься формата или отправляешь несколько команд сразу, я должен назначить тебе штраф.} \\
\texttt{* Если оба игрока наберут больше чем 8 штрафов, вы оба проиграете.} \\
\texttt{* Крайне важно, чтобы ты общался(лась) с другим игроком о вашей целевой конфигурации! *Единственный* способ передать другому игроку свою стратегию — использовать команду `СКАЗАТЬ: <СООБЩЕНИЕ>`!} \\
\\ 
\\ 
\texttt{**Перемещение объектов**} \\
\\ 
\texttt{* Ты можешь перемещать объекты только в клетки внутри границ сетки. Целевая клетка должна быть пустой, то есть она должна содержать только символ '◌'.} \\
\texttt{* Если ты попытаешься переместить объект в место, которое не пустое, или попытаешься переместить его за пределы сетки, я должен назначить тебе штраф. У тебя будет ещё одна попытка.} \\
\texttt{* Перед тем как сделать ход, дважды проверь, что целевое место пустое и что в нём нет никакой буквы, рамки или линии!} \\
\\ 
\\ 
\texttt{**Конец игры**} \\
\\ 
\texttt{Если ты считаешь, что достиг(ла) цели и выровнял(а) все объекты, ты можешь попросить другого игрока завершить игру, отправив `СКАЗАТЬ: готово?`. Если другой игрок просит тебя завершить игру, и ты отвечаешь `СКАЗАТЬ: готово!`, игра завершится.} \\
\\ 
\texttt{Оба игрока выигрывают, если игра закончится в пределах 12 раундов, где один раунд определяется как то, что оба игрока отправили по одной корректной команде.} \\
\\ 
\texttt{**Подсчёт очков:**} \\
\\ 
\texttt{Чем ближе одинаковые объекты расположены на обоих игровых полях, тем больше очков ты получаешь. Штрафы уменьшают твои очки. Сможешь побить рекорд?} \\
\\ 
\texttt{Другой игрок начал игру, отправив это сообщение:} \\
\\ 
\texttt{"Привет! Давай договоримся о целевой конфигурации. Какие у тебя координаты объектов C, L, P на твоём поле? У меня: C в (1,2), L в (1,1), P в (7,7)."} \\
\\ 
\texttt{Какая твоя первая команда?} \\
            }
        }
    }
    & \\ \\

    \theutterance \stepcounter{utterance}  
    & & & \multicolumn{4}{p{0.6\linewidth}}{
        \cellcolor[rgb]{0.9,0.9,0.9}{
            \makecell[{{p{\linewidth}}}]{
                \texttt{\tiny{[GM$\langle$P2]}}
                \texttt{СКАЗАТЬ: Привет! У меня координаты объектов: C на (6,5), L на (5,2), P на (3,3). Давай решать, какую общую конфигурацию выбрать.} \\
            }
        }
    }
     \\ \\

    \theutterance \stepcounter{utterance}  
    & & & \multicolumn{2}{p{0.3\linewidth}}{
        \cellcolor[rgb]{0.9,0.9,0.9}{
            \makecell[{{p{\linewidth}}}]{
                \texttt{\tiny{[GM$|$GM]}}
                \texttt{СКАЗАТЬ: Привет! У меня координаты объектов: C на (6,5), L на (5,2), P на (3,3). Давай решать, какую общую конфигурацию выбрать.} \\
            }
        }
    }
    & & \\ \\

    \theutterance \stepcounter{utterance}  
    & & \multicolumn{4}{p{0.6\linewidth}}{
        \cellcolor[rgb]{0.9,0.9,0.9}{
            \makecell[{{p{\linewidth}}}]{
                \texttt{\tiny{[P1$\langle$GM]}}
                \texttt{Ваше сообщение было передано другому игроку.} \\
\\ 
\\ 
\texttt{Сейчас вы играете раунд 1 из максимум 12.} \\
\texttt{Вы вместе накопили 0 из 8 штрафов.} \\
\texttt{Другой игрок отправил это сообщение:} \\
\\ 
\\ 
\texttt{"Привет! У меня координаты объектов: C на (6,5), L на (5,2), P на (3,3). Давай решать, какую общую конфигурацию выбрать."} \\
\\ 
\\ 
\texttt{Какая твоя следующая команда?} \\
            }
        }
    }
    & \\ \\

    \theutterance \stepcounter{utterance}  
    & \multicolumn{4}{p{0.6\linewidth}}{
        \cellcolor[rgb]{0.9,0.9,0.9}{
            \makecell[{{p{\linewidth}}}]{
                \texttt{\tiny{[P1$\rangle$GM]}}
                \texttt{СКАЗАТЬ: Нам нужно выбрать общую конфигурацию. Предлагаю использовать координаты, которые не заняты рамками и линиями. Посмотрим на свободные клетки. На моём поле свободны (2,1), (3,1), (4,1), (5,1), (6,1), (7,1), (3,2), (4,2), (5,2), (6,2), (7,2), (2,3), (4,3), (5,3), (6,3), (7,3), (1,4), (2,4), (3,4), (4,4), (5,4), (6,4), (7,4), (1,5), (2,5), (5,5), (6,5), (7,5), (1,6), (2,6), (3,6), (4,6), (7,6), (1,7), (2,7), (3,7), (4,7), (6,7). На твоём поле, судя по твоим координатам, рамки могут отличаться. Какие клетки у тебя свободны? Давай выберем три координаты, куда поставим C, L, P, так чтобы они были свободны у обоих.} \\
            }
        }
    }
    & & \\ \\

    \theutterance \stepcounter{utterance}  
    & & & \multicolumn{2}{p{0.3\linewidth}}{
        \cellcolor[rgb]{0.9,0.9,0.9}{
            \makecell[{{p{\linewidth}}}]{
                \texttt{\tiny{[GM$|$GM]}}
                \texttt{СКАЗАТЬ: Нам нужно выбрать общую конфигурацию. Предлагаю использовать координаты, которые не заняты рамками и линиями. Посмотрим на свободные клетки. На моём поле свободны (2,1), (3,1), (4,1), (5,1), (6,1), (7,1), (3,2), (4,2), (5,2), (6,2), (7,2), (2,3), (4,3), (5,3), (6,3), (7,3), (1,4), (2,4), (3,4), (4,4), (5,4), (6,4), (7,4), (1,5), (2,5), (5,5), (6,5), (7,5), (1,6), (2,6), (3,6), (4,6), (7,6), (1,7), (2,7), (3,7), (4,7), (6,7). На твоём поле, судя по твоим координатам, рамки могут отличаться. Какие клетки у тебя свободны? Давай выберем три координаты, куда поставим C, L, P, так чтобы они были свободны у обоих.} \\
            }
        }
    }
    & & \\ \\

    \theutterance \stepcounter{utterance}  
    & & \multicolumn{4}{p{0.6\linewidth}}{
        \cellcolor[rgb]{0.9,0.9,0.9}{
            \makecell[{{p{\linewidth}}}]{
                \texttt{\tiny{[GM$\rangle$P2]}}
                \texttt{Ваше сообщение было передано другому игроку.} \\
\\ 
\\ 
\texttt{Сейчас вы играете раунд 2 из максимум 12.} \\
\texttt{Вы вместе накопили 0 из 8 штрафов.} \\
\texttt{Другой игрок отправил это сообщение:} \\
\\ 
\\ 
\texttt{"Нам нужно выбрать общую конфигурацию. Предлагаю использовать координаты, которые не заняты рамками и линиями. Посмотрим на свободные клетки. На моём поле свободны (2,1), (3,1), (4,1), (5,1), (6,1), (7,1), (3,2), (4,2), (5,2), (6,2), (7,2), (2,3), (4,3), (5,3), (6,3), (7,3), (1,4), (2,4), (3,4), (4,4), (5,4), (6,4), (7,4), (1,5), (2,5), (5,5), (6,5), (7,5), (1,6), (2,6), (3,6), (4,6), (7,6), (1,7), (2,7), (3,7), (4,7), (6,7). На твоём поле, судя по твоим координатам, рамки могут отличаться. Какие клетки у тебя свободны? Давай выберем три координаты, куда поставим C, L, P, так чтобы они были свободны у обоих."} \\
\\ 
\\ 
\texttt{Какая твоя следующая команда?} \\
            }
        }
    }
    & \\ \\

    \theutterance \stepcounter{utterance}  
    & & & \multicolumn{4}{p{0.6\linewidth}}{
        \cellcolor[rgb]{0.9,0.9,0.9}{
            \makecell[{{p{\linewidth}}}]{
                \texttt{\tiny{[GM$\langle$P2]}}
                \texttt{Давай сначала проанализирую своё игровое поле.} \\
\\ 
\texttt{Моё поле (судя по сетке):} \\
\\ 
\texttt{{-} Строка 1: столбцы 1,2,3,4,5,6,7 — все '◌', кроме рамок, но вижу: 1 и 2 разделены линией, 3{-}7 свободны} \\
\texttt{{-} Строка 2: 1,2 свободны; 3,4 разделены; 5{-}L, 6,7 свободны} \\
\texttt{{-} Строка 3: рамка слева, 3{-}P, рамки} \\
\texttt{{-} Строка 4: 1 свободна; 2,3,4,5,6,7 — надо проверить} \\
\texttt{{-} Строка 5: 1 свободна; дальше разрывы; 6{-}C, 7 свободна} \\
\texttt{{-} Строка 6: 1,2 свободны; 3,4,5 свободны; 6 рамка; 7 свободна} \\
\texttt{{-} Строка 7: 1,2 свободны; 3,4,5 свободны; 6,7 свободны} \\
\\ 
\texttt{Посмотрю внимательно на сетку и перечислю все свободные клетки:} \\
\\ 
\texttt{Строка 1: (3,1), (4,1), (5,1), (6,1), (7,1) — (1,1) и (2,1) под вопросом из{-}за символов, но судя по сетке там '◌', значит свободны. Перепроверю: (1,1)='◌', (2,1)='◌'. Значит все 1{-}7 свободны, кроме разделителей? Но в правилах сказано: "Целевая клетка должна быть пустой, то есть она должна содержать только символ '◌'." На сетке в строке 1: позиции 1 и 2 показаны как '◌', а 3{-}7 как '◌◌◌'. Так что (1,1), (2,1), (3,1), (4,1), (5,1), (6,1), (7,1) свободны.} \\
\\ 
\texttt{Строка 2: (1,2)='◌', (2,2)='◌', (3,2) — там '┤'? Нет, на сетке: "║◌├─┤◌L◌║" — позиции: 1='◌', 2='├', 3='─', 4='┤', 5='◌', 6='L', 7='◌'. Свободны: (1,2), (5,2), (7,2).} \\
\\ 
\texttt{Строка 3: "╟─┤P├─┬─╢" — позиции: 1='╟', 2='─', 3='┤', 4='P', 5='├', 6='─', 7='┬'. Свободных нет, кроме P.} \\
\\ 
\texttt{Строка 4: "║◌│◌│◌│◌║" — позиции: 1='◌', 2='│', 3='◌', 4='│', 5='◌', 6='│', 7='◌'. Свободны: (1,4), (3,4), (5,4), (7,4).} \\
\\ 
\texttt{Строка 5: "╟─┼─┘C│◌║" — позиции: 1='╟', 2='─', 3='┼', 4='─', 5='┘', 6='C', 7='│'? Нет: "╟─┼─┘C│◌║" — это 8 символов? Считаю: ╟ (1), ─ (2), ┼ (3), ─ (4), ┘ (5), C (6), │ (7), ◌ (8)? Но x{-}координаты 1{-}7. Посмотрю: "╟─┼─┘C│◌║" — первые 7: 1=╟, 2=─, 3=┼, 4=─, 5=┘, 6=C, 7=│, а потом ◌║ для строки 5? Но в сетке строка 5: "╟─┼─┘C│◌║" и строка 6: "║◌│◌◌┌┴─╢". Возможно, в строке 5 после C идёт │ и ◌, но тогда x=7 это │? Перечитаю: "╟─┼─┘C│◌║" — если это строка, то символы: 1:╟, 2:─, 3:┼, 4:─, 5:┘, 6:C, 7:│, и затем ║ — это конец строки? Но тогда где ◌? Может быть, "╟─┼─┘C│◌║" — ◌ это часть строки 5? Всего символов: ╟─┼─┘C│◌║ — 9 штук. Но x от 1 до 7, значит ◌ на x=7? Тогда: 1:╟, 2:─, 3:┼, 4:─, 5:┘, 6:C, 7:◌, и ║ это граница. Да, правильно. Тогда строка 5: свободна только (7,5)='◌'.} \\
\\ 
\texttt{Строка 6: "║◌│◌◌┌┴─╢" — позиции: 1='◌', 2='│', 3='◌', 4='◌', 5='┌', 6='┴', 7='─'. Свободны: (1,6), (3,6), (4,6). Но (7,6) не свободна — там '─'.} \\
\\ 
\texttt{Строка 7: "║◌│◌◌│◌◌║" — позиции: 1='◌', 2='│', 3='◌', 4='◌', 5='│', 6='◌', 7='◌'. Свободны: (1,7), (3,7), (4,7), (6,7), (7,7).} \\
\\ 
\texttt{Итак, мои свободные клетки: (1,1){-}(7,1), (1,2), (5,2), (7,2), (1,4), (3,4), (5,4), (7,4), (7,5), (1,6), (3,6), (4,6), (1,7), (3,7), (4,7), (6,7), (7,7).} \\
\\ 
\texttt{Сравню со списком другого игрока. У него свободны многие клетки, включая (2,1), (3,1), (4,1), (5,1), (6,1), (7,1), а также (3,2), (4,2), (5,2), (6,2), (7,2) — у меня (3,2) нет. Также у него есть (2,3), (4,3), (5,3), (6,3), (7,3) — у меня в строке 3 почти всё занято. У него (1,4){-}(7,4) все свободны — у меня только (1,4),(3,4),(5,4),(7,4). У него (1,5),(2,5),(5,5),(6,5),(7,5) — у меня только (7,5). У него (1,6){-}(4,6),(7,6) — у меня (1,6),(3,6),(4,6). У него (1,7){-}(4,7),(6,7) — у меня (1,7),(3,7),(4,7),(6,7),(7,7).} \\
\\ 
\texttt{Общие свободные клетки, которые есть у обоих (предположительно, судя по его списку):} \\
\texttt{Строка 1: (1,1){-}(7,1), у него (2,1){-}(7,1) — (1,1) он не упомянул; проверим: "На моём поле свободны (2,1), (3,1), (4,1), (5,1), (6,1), (7,1)" — (1,1) у него, видимо, занята. Значит общие: (2,1){-}(7,1).} \\
\texttt{Строка 2:} \\
            }
        }
    }
     \\ \\

    \theutterance \stepcounter{utterance}  
    & & & \multicolumn{2}{p{0.3\linewidth}}{
        \cellcolor[rgb]{0.9,0.9,0.9}{
            \makecell[{{p{\linewidth}}}]{
                \texttt{\tiny{[GM$|$GM]}}
                \texttt{Давай сначала проанализирую своё игровое поле.} \\
\\ 
\texttt{Моё поле (судя по сетке):} \\
\\ 
\texttt{{-} Строка 1: столбцы 1,2,3,4,5,6,7 — все '◌', кроме рамок, но вижу: 1 и 2 разделены линией, 3{-}7 свободны} \\
\texttt{{-} Строка 2: 1,2 свободны; 3,4 разделены; 5{-}L, 6,7 свободны} \\
\texttt{{-} Строка 3: рамка слева, 3{-}P, рамки} \\
\texttt{{-} Строка 4: 1 свободна; 2,3,4,5,6,7 — надо проверить} \\
\texttt{{-} Строка 5: 1 свободна; дальше разрывы; 6{-}C, 7 свободна} \\
\texttt{{-} Строка 6: 1,2 свободны; 3,4,5 свободны; 6 рамка; 7 свободна} \\
\texttt{{-} Строка 7: 1,2 свободны; 3,4,5 свободны; 6,7 свободны} \\
\\ 
\texttt{Посмотрю внимательно на сетку и перечислю все свободные клетки:} \\
\\ 
\texttt{Строка 1: (3,1), (4,1), (5,1), (6,1), (7,1) — (1,1) и (2,1) под вопросом из{-}за символов, но судя по сетке там '◌', значит свободны. Перепроверю: (1,1)='◌', (2,1)='◌'. Значит все 1{-}7 свободны, кроме разделителей? Но в правилах сказано: "Целевая клетка должна быть пустой, то есть она должна содержать только символ '◌'." На сетке в строке 1: позиции 1 и 2 показаны как '◌', а 3{-}7 как '◌◌◌'. Так что (1,1), (2,1), (3,1), (4,1), (5,1), (6,1), (7,1) свободны.} \\
\\ 
\texttt{Строка 2: (1,2)='◌', (2,2)='◌', (3,2) — там '┤'? Нет, на сетке: "║◌├─┤◌L◌║" — позиции: 1='◌', 2='├', 3='─', 4='┤', 5='◌', 6='L', 7='◌'. Свободны: (1,2), (5,2), (7,2).} \\
\\ 
\texttt{Строка 3: "╟─┤P├─┬─╢" — позиции: 1='╟', 2='─', 3='┤', 4='P', 5='├', 6='─', 7='┬'. Свободных нет, кроме P.} \\
\\ 
\texttt{Строка 4: "║◌│◌│◌│◌║" — позиции: 1='◌', 2='│', 3='◌', 4='│', 5='◌', 6='│', 7='◌'. Свободны: (1,4), (3,4), (5,4), (7,4).} \\
\\ 
\texttt{Строка 5: "╟─┼─┘C│◌║" — позиции: 1='╟', 2='─', 3='┼', 4='─', 5='┘', 6='C', 7='│'? Нет: "╟─┼─┘C│◌║" — это 8 символов? Считаю: ╟ (1), ─ (2), ┼ (3), ─ (4), ┘ (5), C (6), │ (7), ◌ (8)? Но x{-}координаты 1{-}7. Посмотрю: "╟─┼─┘C│◌║" — первые 7: 1=╟, 2=─, 3=┼, 4=─, 5=┘, 6=C, 7=│, а потом ◌║ для строки 5? Но в сетке строка 5: "╟─┼─┘C│◌║" и строка 6: "║◌│◌◌┌┴─╢". Возможно, в строке 5 после C идёт │ и ◌, но тогда x=7 это │? Перечитаю: "╟─┼─┘C│◌║" — если это строка, то символы: 1:╟, 2:─, 3:┼, 4:─, 5:┘, 6:C, 7:│, и затем ║ — это конец строки? Но тогда где ◌? Может быть, "╟─┼─┘C│◌║" — ◌ это часть строки 5? Всего символов: ╟─┼─┘C│◌║ — 9 штук. Но x от 1 до 7, значит ◌ на x=7? Тогда: 1:╟, 2:─, 3:┼, 4:─, 5:┘, 6:C, 7:◌, и ║ это граница. Да, правильно. Тогда строка 5: свободна только (7,5)='◌'.} \\
\\ 
\texttt{Строка 6: "║◌│◌◌┌┴─╢" — позиции: 1='◌', 2='│', 3='◌', 4='◌', 5='┌', 6='┴', 7='─'. Свободны: (1,6), (3,6), (4,6). Но (7,6) не свободна — там '─'.} \\
\\ 
\texttt{Строка 7: "║◌│◌◌│◌◌║" — позиции: 1='◌', 2='│', 3='◌', 4='◌', 5='│', 6='◌', 7='◌'. Свободны: (1,7), (3,7), (4,7), (6,7), (7,7).} \\
\\ 
\texttt{Итак, мои свободные клетки: (1,1){-}(7,1), (1,2), (5,2), (7,2), (1,4), (3,4), (5,4), (7,4), (7,5), (1,6), (3,6), (4,6), (1,7), (3,7), (4,7), (6,7), (7,7).} \\
\\ 
\texttt{Сравню со списком другого игрока. У него свободны многие клетки, включая (2,1), (3,1), (4,1), (5,1), (6,1), (7,1), а также (3,2), (4,2), (5,2), (6,2), (7,2) — у меня (3,2) нет. Также у него есть (2,3), (4,3), (5,3), (6,3), (7,3) — у меня в строке 3 почти всё занято. У него (1,4){-}(7,4) все свободны — у меня только (1,4),(3,4),(5,4),(7,4). У него (1,5),(2,5),(5,5),(6,5),(7,5) — у меня только (7,5). У него (1,6){-}(4,6),(7,6) — у меня (1,6),(3,6),(4,6). У него (1,7){-}(4,7),(6,7) — у меня (1,7),(3,7),(4,7),(6,7),(7,7).} \\
\\ 
\texttt{Общие свободные клетки, которые есть у обоих (предположительно, судя по его списку):} \\
\texttt{Строка 1: (1,1){-}(7,1), у него (2,1){-}(7,1) — (1,1) он не упомянул; проверим: "На моём поле свободны (2,1), (3,1), (4,1), (5,1), (6,1), (7,1)" — (1,1) у него, видимо, занята. Значит общие: (2,1){-}(7,1).} \\
\texttt{Строка 2:} \\
            }
        }
    }
    & & \\ \\

    \theutterance \stepcounter{utterance}  
    & & & \multicolumn{2}{p{0.3\linewidth}}{
        \cellcolor[rgb]{0.9,0.9,0.9}{
            \makecell[{{p{\linewidth}}}]{
                \texttt{\tiny{[GM$|$GM]}}
                \texttt{Invalid response format} \\
            }
        }
    }
    & & \\ \\

    \theutterance \stepcounter{utterance}  
    & & \multicolumn{4}{p{0.6\linewidth}}{
        \cellcolor[rgb]{0.9,0.9,0.9}{
            \makecell[{{p{\linewidth}}}]{
                \texttt{\tiny{[GM$\rangle$P2]}}
                \texttt{Штраф: Ваше сообщение не соответствует ожидаемому формату!} \\
\\ 
\texttt{Убедитесь, что ваш ответ содержит только `СКАЗАТЬ: <СООБЩЕНИЕ>` или `ДВИГАТЬСЯ: <ОБЪЕКТ>, (<X>, <Y>)` и ничего больше!} \\
\\ 
\texttt{Вы вместе накопили 1 из 8 штрафов.} \\
\texttt{Пожалуйста, попробуйте ещё раз!} \\
            }
        }
    }
    & \\ \\

    \theutterance \stepcounter{utterance}  
    & & & \multicolumn{4}{p{0.6\linewidth}}{
        \cellcolor[rgb]{0.9,0.9,0.9}{
            \makecell[{{p{\linewidth}}}]{
                \texttt{\tiny{[GM$\langle$P2]}}
                \texttt{СКАЗАТЬ: Понял, извиняюсь за длинное сообщение. У меня свободные клетки: (2,1){-}(7,1), (5,2), (7,2), (1,4), (3,4), (5,4), (7,4), (7,5), (1,6), (3,6), (4,6), (1,7), (3,7), (4,7), (6,7), (7,7). Общие с твоими: (2,1){-}(7,1), (5,2), (7,2), (5,4), (7,4), (7,5), (1,6), (3,6), (4,6), (1,7), (3,7), (4,7), (6,7). Предлагаю выбрать три из них для C, L, P.} \\
            }
        }
    }
     \\ \\

    \theutterance \stepcounter{utterance}  
    & & & \multicolumn{2}{p{0.3\linewidth}}{
        \cellcolor[rgb]{0.9,0.9,0.9}{
            \makecell[{{p{\linewidth}}}]{
                \texttt{\tiny{[GM$|$GM]}}
                \texttt{СКАЗАТЬ: Понял, извиняюсь за длинное сообщение. У меня свободные клетки: (2,1){-}(7,1), (5,2), (7,2), (1,4), (3,4), (5,4), (7,4), (7,5), (1,6), (3,6), (4,6), (1,7), (3,7), (4,7), (6,7), (7,7). Общие с твоими: (2,1){-}(7,1), (5,2), (7,2), (5,4), (7,4), (7,5), (1,6), (3,6), (4,6), (1,7), (3,7), (4,7), (6,7). Предлагаю выбрать три из них для C, L, P.} \\
            }
        }
    }
    & & \\ \\

    \theutterance \stepcounter{utterance}  
    & & \multicolumn{4}{p{0.6\linewidth}}{
        \cellcolor[rgb]{0.9,0.9,0.9}{
            \makecell[{{p{\linewidth}}}]{
                \texttt{\tiny{[P1$\langle$GM]}}
                \texttt{Ваше сообщение было передано другому игроку.} \\
\\ 
\\ 
\texttt{Сейчас вы играете раунд 2 из максимум 12.} \\
\texttt{Вы вместе накопили 1 из 8 штрафов.} \\
\texttt{Другой игрок отправил это сообщение:} \\
\\ 
\\ 
\texttt{"Понял, извиняюсь за длинное сообщение. У меня свободные клетки: (2,1){-}(7,1), (5,2), (7,2), (1,4), (3,4), (5,4), (7,4), (7,5), (1,6), (3,6), (4,6), (1,7), (3,7), (4,7), (6,7), (7,7). Общие с твоими: (2,1){-}(7,1), (5,2), (7,2), (5,4), (7,4), (7,5), (1,6), (3,6), (4,6), (1,7), (3,7), (4,7), (6,7). Предлагаю выбрать три из них для C, L, P."} \\
\\ 
\\ 
\texttt{Какая твоя следующая команда?} \\
            }
        }
    }
    & \\ \\

    \theutterance \stepcounter{utterance}  
    & \multicolumn{4}{p{0.6\linewidth}}{
        \cellcolor[rgb]{0.9,0.9,0.9}{
            \makecell[{{p{\linewidth}}}]{
                \texttt{\tiny{[P1$\rangle$GM]}}
                \texttt{СКАЗАТЬ: Отлично! Из общих свободных клеток предлагаю такую конфигурацию: C на (5,4), L на (3,6), P на (7,2). Проверь, пожалуйста, свободны ли они у тебя. Если да, то начнём двигать объекты. Как тебе такой план?} \\
            }
        }
    }
    & & \\ \\

    \theutterance \stepcounter{utterance}  
    & & & \multicolumn{2}{p{0.3\linewidth}}{
        \cellcolor[rgb]{0.9,0.9,0.9}{
            \makecell[{{p{\linewidth}}}]{
                \texttt{\tiny{[GM$|$GM]}}
                \texttt{СКАЗАТЬ: Отлично! Из общих свободных клеток предлагаю такую конфигурацию: C на (5,4), L на (3,6), P на (7,2). Проверь, пожалуйста, свободны ли они у тебя. Если да, то начнём двигать объекты. Как тебе такой план?} \\
            }
        }
    }
    & & \\ \\

    \theutterance \stepcounter{utterance}  
    & & \multicolumn{4}{p{0.6\linewidth}}{
        \cellcolor[rgb]{0.9,0.9,0.9}{
            \makecell[{{p{\linewidth}}}]{
                \texttt{\tiny{[GM$\rangle$P2]}}
                \texttt{Ваше сообщение было передано другому игроку.} \\
\\ 
\\ 
\texttt{Сейчас вы играете раунд 3 из максимум 12.} \\
\texttt{Вы вместе накопили 1 из 8 штрафов.} \\
\texttt{Другой игрок отправил это сообщение:} \\
\\ 
\\ 
\texttt{"Отлично! Из общих свободных клеток предлагаю такую конфигурацию: C на (5,4), L на (3,6), P на (7,2). Проверь, пожалуйста, свободны ли они у тебя. Если да, то начнём двигать объекты. Как тебе такой план?"} \\
\\ 
\\ 
\texttt{Какая твоя следующая команда?} \\
            }
        }
    }
    & \\ \\

    \theutterance \stepcounter{utterance}  
    & & & \multicolumn{4}{p{0.6\linewidth}}{
        \cellcolor[rgb]{0.9,0.9,0.9}{
            \makecell[{{p{\linewidth}}}]{
                \texttt{\tiny{[GM$\langle$P2]}}
                \texttt{Проверю предложенные координаты на своём поле: (5,4) — у меня там '◌', свободна; (3,6) — '◌', свободна; (7,2) — '◌', свободна. Все три клетки свободны, отлично.} \\
\\ 
\texttt{СКАЗАТЬ: Отличный план! Все три клетки у меня свободны — (5,4), (3,6), (7,2). Начинаем двигать. Я перемещу P на (7,2) первым ходом.} \\
            }
        }
    }
     \\ \\

    \theutterance \stepcounter{utterance}  
    & & & \multicolumn{2}{p{0.3\linewidth}}{
        \cellcolor[rgb]{0.9,0.9,0.9}{
            \makecell[{{p{\linewidth}}}]{
                \texttt{\tiny{[GM$|$GM]}}
                \texttt{Проверю предложенные координаты на своём поле: (5,4) — у меня там '◌', свободна; (3,6) — '◌', свободна; (7,2) — '◌', свободна. Все три клетки свободны, отлично.} \\
\\ 
\texttt{СКАЗАТЬ: Отличный план! Все три клетки у меня свободны — (5,4), (3,6), (7,2). Начинаем двигать. Я перемещу P на (7,2) первым ходом.} \\
            }
        }
    }
    & & \\ \\

    \theutterance \stepcounter{utterance}  
    & & & \multicolumn{2}{p{0.3\linewidth}}{
        \cellcolor[rgb]{0.9,0.9,0.9}{
            \makecell[{{p{\linewidth}}}]{
                \texttt{\tiny{[GM$|$GM]}}
                \texttt{Invalid format: head is not empty: 'Проверю предложенные координаты на своём поле: (5,4) — у меня там '◌', свободна; (3,6) — '◌', свободна; (7,2) — '◌', свободна. Все три клетки свободны, отлично.} \\
\\ 
\texttt{'} \\
            }
        }
    }
    & & \\ \\

    \theutterance \stepcounter{utterance}  
    & & \multicolumn{4}{p{0.6\linewidth}}{
        \cellcolor[rgb]{0.9,0.9,0.9}{
            \makecell[{{p{\linewidth}}}]{
                \texttt{\tiny{[GM$\rangle$P2]}}
                \texttt{Штраф: Ваше сообщение не должно содержать ничего перед командой!} \\
\\ 
\texttt{Убедитесь, что ваш ответ содержит только `СКАЗАТЬ: <СООБЩЕНИЕ>` или `ДВИГАТЬСЯ: <ОБЪЕКТ>, (<X>, <Y>)` и ничего больше!} \\
\\ 
\texttt{Вы вместе накопили 2 из 8 штрафов.} \\
\texttt{Пожалуйста, попробуйте ещё раз!} \\
            }
        }
    }
    & \\ \\

    \theutterance \stepcounter{utterance}  
    & & & \multicolumn{4}{p{0.6\linewidth}}{
        \cellcolor[rgb]{0.9,0.9,0.9}{
            \makecell[{{p{\linewidth}}}]{
                \texttt{\tiny{[GM$\langle$P2]}}
                \texttt{СКАЗАТЬ: Все три клетки у меня свободны — (5,4), (3,6), (7,2). Начинаем двигать. Я перемещу P на (7,2) первым.} \\
            }
        }
    }
     \\ \\

    \theutterance \stepcounter{utterance}  
    & & & \multicolumn{2}{p{0.3\linewidth}}{
        \cellcolor[rgb]{0.9,0.9,0.9}{
            \makecell[{{p{\linewidth}}}]{
                \texttt{\tiny{[GM$|$GM]}}
                \texttt{СКАЗАТЬ: Все три клетки у меня свободны — (5,4), (3,6), (7,2). Начинаем двигать. Я перемещу P на (7,2) первым.} \\
            }
        }
    }
    & & \\ \\

    \theutterance \stepcounter{utterance}  
    & & \multicolumn{4}{p{0.6\linewidth}}{
        \cellcolor[rgb]{0.9,0.9,0.9}{
            \makecell[{{p{\linewidth}}}]{
                \texttt{\tiny{[P1$\langle$GM]}}
                \texttt{Ваше сообщение было передано другому игроку.} \\
\\ 
\\ 
\texttt{Сейчас вы играете раунд 3 из максимум 12.} \\
\texttt{Вы вместе накопили 2 из 8 штрафов.} \\
\texttt{Другой игрок отправил это сообщение:} \\
\\ 
\\ 
\texttt{"Все три клетки у меня свободны — (5,4), (3,6), (7,2). Начинаем двигать. Я перемещу P на (7,2) первым."} \\
\\ 
\\ 
\texttt{Какая твоя следующая команда?} \\
            }
        }
    }
    & \\ \\

    \theutterance \stepcounter{utterance}  
    & \multicolumn{4}{p{0.6\linewidth}}{
        \cellcolor[rgb]{0.9,0.9,0.9}{
            \makecell[{{p{\linewidth}}}]{
                \texttt{\tiny{[P1$\rangle$GM]}}
                \texttt{ДВИГАТЬСЯ: P, (7, 2)} \\
            }
        }
    }
    & & \\ \\

    \theutterance \stepcounter{utterance}  
    & & & \multicolumn{2}{p{0.3\linewidth}}{
        \cellcolor[rgb]{0.9,0.9,0.9}{
            \makecell[{{p{\linewidth}}}]{
                \texttt{\tiny{[GM$|$GM]}}
                \texttt{ДВИГАТЬСЯ: P, (7, 2)} \\
            }
        }
    }
    & & \\ \\

    \theutterance \stepcounter{utterance}  
    & & & \multicolumn{2}{p{0.3\linewidth}}{
        \cellcolor[rgb]{0.9,0.9,0.9}{
            \makecell[{{p{\linewidth}}}]{
                \texttt{\tiny{[GM$|$GM]}}
                \texttt{P moved to (7, 2)} \\
            }
        }
    }
    & & \\ \\

    \theutterance \stepcounter{utterance}  
    & & \multicolumn{4}{p{0.6\linewidth}}{
        \cellcolor[rgb]{0.9,0.9,0.9}{
            \makecell[{{p{\linewidth}}}]{
                \texttt{\tiny{[GM$\rangle$P2]}}
                \texttt{Ваше сообщение было передано другому игроку.} \\
\\ 
\\ 
\texttt{Сейчас вы играете раунд 4 из максимум 12.} \\
\texttt{Вы вместе накопили 2 из 8 штрафов.} \\
\texttt{Другой игрок переместил объект на своей сетке.} \\
\\ 
\\ 
\texttt{Какая твоя следующая команда?} \\
            }
        }
    }
    & \\ \\

    \theutterance \stepcounter{utterance}  
    & & & \multicolumn{4}{p{0.6\linewidth}}{
        \cellcolor[rgb]{0.9,0.9,0.9}{
            \makecell[{{p{\linewidth}}}]{
                \texttt{\tiny{[GM$\langle$P2]}}
                \texttt{ДВИГАТЬСЯ: P, (7,2)} \\
            }
        }
    }
     \\ \\

    \theutterance \stepcounter{utterance}  
    & & & \multicolumn{2}{p{0.3\linewidth}}{
        \cellcolor[rgb]{0.9,0.9,0.9}{
            \makecell[{{p{\linewidth}}}]{
                \texttt{\tiny{[GM$|$GM]}}
                \texttt{ДВИГАТЬСЯ: P, (7,2)} \\
            }
        }
    }
    & & \\ \\

    \theutterance \stepcounter{utterance}  
    & & & \multicolumn{2}{p{0.3\linewidth}}{
        \cellcolor[rgb]{0.9,0.9,0.9}{
            \makecell[{{p{\linewidth}}}]{
                \texttt{\tiny{[GM$|$GM]}}
                \texttt{P moved to (7, 2)} \\
            }
        }
    }
    & & \\ \\

    \theutterance \stepcounter{utterance}  
    & & \multicolumn{4}{p{0.6\linewidth}}{
        \cellcolor[rgb]{0.9,0.9,0.9}{
            \makecell[{{p{\linewidth}}}]{
                \texttt{\tiny{[P1$\langle$GM]}}
                \texttt{Успешно перемещён объект "P" в (7,2). Ваша обновлённая сетка выглядит так:} \\
\texttt{```} \\
\texttt{ 1234567} \\
\texttt{╔═╤═╤═══╗} \\
\texttt{║L│◌│◌◌◌║ 1} \\
\texttt{║C├─┤◌◌P║ 2} \\
\texttt{╟─┤◌├─┬─╢ 3} \\
\texttt{║◌│◌│◌│◌║ 4} \\
\texttt{╟─┼─┘◌│◌║ 5} \\
\texttt{║◌│◌◌┌┴─╢ 6} \\
\texttt{║◌│◌◌│◌◌║ 7} \\
\texttt{╚═╧══╧══╝} \\
\texttt{```} \\
\texttt{Сейчас вы играете раунд 4 из максимум 12.} \\
\texttt{Вы вместе накопили 2 из 8 штрафов.} \\
\texttt{Другой игрок переместил объект на своей сетке.} \\
\\ 
\\ 
\texttt{Какая твоя следующая команда?} \\
            }
        }
    }
    & \\ \\

    \theutterance \stepcounter{utterance}  
    & \multicolumn{4}{p{0.6\linewidth}}{
        \cellcolor[rgb]{0.9,0.9,0.9}{
            \makecell[{{p{\linewidth}}}]{
                \texttt{\tiny{[P1$\rangle$GM]}}
                \texttt{СКАЗАТЬ: Я передвинул P на (7,2). Теперь двигаем L на (3,6). Ты тоже можешь двигать L на (3,6), если у тебя там свободно.} \\
            }
        }
    }
    & & \\ \\

    \theutterance \stepcounter{utterance}  
    & & & \multicolumn{2}{p{0.3\linewidth}}{
        \cellcolor[rgb]{0.9,0.9,0.9}{
            \makecell[{{p{\linewidth}}}]{
                \texttt{\tiny{[GM$|$GM]}}
                \texttt{СКАЗАТЬ: Я передвинул P на (7,2). Теперь двигаем L на (3,6). Ты тоже можешь двигать L на (3,6), если у тебя там свободно.} \\
            }
        }
    }
    & & \\ \\

    \theutterance \stepcounter{utterance}  
    & & \multicolumn{4}{p{0.6\linewidth}}{
        \cellcolor[rgb]{0.9,0.9,0.9}{
            \makecell[{{p{\linewidth}}}]{
                \texttt{\tiny{[GM$\rangle$P2]}}
                \texttt{Успешно перемещён объект "P" в (7,2). Ваша обновлённая сетка выглядит так:} \\
\texttt{```} \\
\texttt{ 1234567} \\
\texttt{╔═╤═╤═══╗} \\
\texttt{║◌│◌│◌◌◌║ 1} \\
\texttt{║◌├─┤◌LP║ 2} \\
\texttt{╟─┤◌├─┬─╢ 3} \\
\texttt{║◌│◌│◌│◌║ 4} \\
\texttt{╟─┼─┘C│◌║ 5} \\
\texttt{║◌│◌◌┌┴─╢ 6} \\
\texttt{║◌│◌◌│◌◌║ 7} \\
\texttt{╚═╧══╧══╝} \\
\texttt{```} \\
\texttt{Сейчас вы играете раунд 5 из максимум 12.} \\
\texttt{Вы вместе накопили 2 из 8 штрафов.} \\
\texttt{Другой игрок отправил это сообщение:} \\
\\ 
\\ 
\texttt{"Я передвинул P на (7,2). Теперь двигаем L на (3,6). Ты тоже можешь двигать L на (3,6), если у тебя там свободно."} \\
\\ 
\\ 
\texttt{Какая твоя следующая команда?} \\
            }
        }
    }
    & \\ \\

    \theutterance \stepcounter{utterance}  
    & & & \multicolumn{4}{p{0.6\linewidth}}{
        \cellcolor[rgb]{0.9,0.9,0.9}{
            \makecell[{{p{\linewidth}}}]{
                \texttt{\tiny{[GM$\langle$P2]}}
                \texttt{ДВИГАТЬСЯ: L, (3,6)} \\
            }
        }
    }
     \\ \\

    \theutterance \stepcounter{utterance}  
    & & & \multicolumn{2}{p{0.3\linewidth}}{
        \cellcolor[rgb]{0.9,0.9,0.9}{
            \makecell[{{p{\linewidth}}}]{
                \texttt{\tiny{[GM$|$GM]}}
                \texttt{ДВИГАТЬСЯ: L, (3,6)} \\
            }
        }
    }
    & & \\ \\

    \theutterance \stepcounter{utterance}  
    & & & \multicolumn{2}{p{0.3\linewidth}}{
        \cellcolor[rgb]{0.9,0.9,0.9}{
            \makecell[{{p{\linewidth}}}]{
                \texttt{\tiny{[GM$|$GM]}}
                \texttt{L moved to (3, 6)} \\
            }
        }
    }
    & & \\ \\

    \theutterance \stepcounter{utterance}  
    & & \multicolumn{4}{p{0.6\linewidth}}{
        \cellcolor[rgb]{0.9,0.9,0.9}{
            \makecell[{{p{\linewidth}}}]{
                \texttt{\tiny{[P1$\langle$GM]}}
                \texttt{Ваше сообщение было передано другому игроку.} \\
\\ 
\\ 
\texttt{Сейчас вы играете раунд 5 из максимум 12.} \\
\texttt{Вы вместе накопили 2 из 8 штрафов.} \\
\texttt{Другой игрок переместил объект на своей сетке.} \\
\\ 
\\ 
\texttt{Какая твоя следующая команда?} \\
            }
        }
    }
    & \\ \\

    \theutterance \stepcounter{utterance}  
    & \multicolumn{4}{p{0.6\linewidth}}{
        \cellcolor[rgb]{0.9,0.9,0.9}{
            \makecell[{{p{\linewidth}}}]{
                \texttt{\tiny{[P1$\rangle$GM]}}
                \texttt{ДВИГАТЬСЯ: L, (3, 6)} \\
            }
        }
    }
    & & \\ \\

    \theutterance \stepcounter{utterance}  
    & & & \multicolumn{2}{p{0.3\linewidth}}{
        \cellcolor[rgb]{0.9,0.9,0.9}{
            \makecell[{{p{\linewidth}}}]{
                \texttt{\tiny{[GM$|$GM]}}
                \texttt{ДВИГАТЬСЯ: L, (3, 6)} \\
            }
        }
    }
    & & \\ \\

    \theutterance \stepcounter{utterance}  
    & & & \multicolumn{2}{p{0.3\linewidth}}{
        \cellcolor[rgb]{0.9,0.9,0.9}{
            \makecell[{{p{\linewidth}}}]{
                \texttt{\tiny{[GM$|$GM]}}
                \texttt{L moved to (3, 6)} \\
            }
        }
    }
    & & \\ \\

    \theutterance \stepcounter{utterance}  
    & & \multicolumn{4}{p{0.6\linewidth}}{
        \cellcolor[rgb]{0.9,0.9,0.9}{
            \makecell[{{p{\linewidth}}}]{
                \texttt{\tiny{[GM$\rangle$P2]}}
                \texttt{Успешно перемещён объект "L" в (3,6). Ваша обновлённая сетка выглядит так:} \\
\texttt{```} \\
\texttt{ 1234567} \\
\texttt{╔═╤═╤═══╗} \\
\texttt{║◌│◌│◌◌◌║ 1} \\
\texttt{║◌├─┤◌◌P║ 2} \\
\texttt{╟─┤◌├─┬─╢ 3} \\
\texttt{║◌│◌│◌│◌║ 4} \\
\texttt{╟─┼─┘C│◌║ 5} \\
\texttt{║◌│L◌┌┴─╢ 6} \\
\texttt{║◌│◌◌│◌◌║ 7} \\
\texttt{╚═╧══╧══╝} \\
\texttt{```} \\
\texttt{Сейчас вы играете раунд 6 из максимум 12.} \\
\texttt{Вы вместе накопили 2 из 8 штрафов.} \\
\texttt{Другой игрок переместил объект на своей сетке.} \\
\\ 
\\ 
\texttt{Какая твоя следующая команда?} \\
            }
        }
    }
    & \\ \\

    \theutterance \stepcounter{utterance}  
    & & & \multicolumn{4}{p{0.6\linewidth}}{
        \cellcolor[rgb]{0.9,0.9,0.9}{
            \makecell[{{p{\linewidth}}}]{
                \texttt{\tiny{[GM$\langle$P2]}}
                \texttt{ДВИГАТЬСЯ: C, (5,4)} \\
            }
        }
    }
     \\ \\

    \theutterance \stepcounter{utterance}  
    & & & \multicolumn{2}{p{0.3\linewidth}}{
        \cellcolor[rgb]{0.9,0.9,0.9}{
            \makecell[{{p{\linewidth}}}]{
                \texttt{\tiny{[GM$|$GM]}}
                \texttt{ДВИГАТЬСЯ: C, (5,4)} \\
            }
        }
    }
    & & \\ \\

    \theutterance \stepcounter{utterance}  
    & & & \multicolumn{2}{p{0.3\linewidth}}{
        \cellcolor[rgb]{0.9,0.9,0.9}{
            \makecell[{{p{\linewidth}}}]{
                \texttt{\tiny{[GM$|$GM]}}
                \texttt{C moved to (5, 4)} \\
            }
        }
    }
    & & \\ \\

    \theutterance \stepcounter{utterance}  
    & & \multicolumn{4}{p{0.6\linewidth}}{
        \cellcolor[rgb]{0.9,0.9,0.9}{
            \makecell[{{p{\linewidth}}}]{
                \texttt{\tiny{[P1$\langle$GM]}}
                \texttt{Успешно перемещён объект "L" в (3,6). Ваша обновлённая сетка выглядит так:} \\
\texttt{```} \\
\texttt{ 1234567} \\
\texttt{╔═╤═╤═══╗} \\
\texttt{║◌│◌│◌◌◌║ 1} \\
\texttt{║C├─┤◌◌P║ 2} \\
\texttt{╟─┤◌├─┬─╢ 3} \\
\texttt{║◌│◌│◌│◌║ 4} \\
\texttt{╟─┼─┘◌│◌║ 5} \\
\texttt{║◌│L◌┌┴─╢ 6} \\
\texttt{║◌│◌◌│◌◌║ 7} \\
\texttt{╚═╧══╧══╝} \\
\texttt{```} \\
\texttt{Сейчас вы играете раунд 6 из максимум 12.} \\
\texttt{Вы вместе накопили 2 из 8 штрафов.} \\
\texttt{Другой игрок переместил объект на своей сетке.} \\
\\ 
\\ 
\texttt{Какая твоя следующая команда?} \\
            }
        }
    }
    & \\ \\

    \theutterance \stepcounter{utterance}  
    & \multicolumn{4}{p{0.6\linewidth}}{
        \cellcolor[rgb]{0.9,0.9,0.9}{
            \makecell[{{p{\linewidth}}}]{
                \texttt{\tiny{[P1$\rangle$GM]}}
                \texttt{СКАЗАТЬ: Я передвинул L на (3,6). Теперь осталось C на (5,4). Двигаем C. У меня там свободно, проверь у себя.} \\
            }
        }
    }
    & & \\ \\

    \theutterance \stepcounter{utterance}  
    & & & \multicolumn{2}{p{0.3\linewidth}}{
        \cellcolor[rgb]{0.9,0.9,0.9}{
            \makecell[{{p{\linewidth}}}]{
                \texttt{\tiny{[GM$|$GM]}}
                \texttt{СКАЗАТЬ: Я передвинул L на (3,6). Теперь осталось C на (5,4). Двигаем C. У меня там свободно, проверь у себя.} \\
            }
        }
    }
    & & \\ \\

    \theutterance \stepcounter{utterance}  
    & & \multicolumn{4}{p{0.6\linewidth}}{
        \cellcolor[rgb]{0.9,0.9,0.9}{
            \makecell[{{p{\linewidth}}}]{
                \texttt{\tiny{[GM$\rangle$P2]}}
                \texttt{Успешно перемещён объект "C" в (5,4). Ваша обновлённая сетка выглядит так:} \\
\texttt{```} \\
\texttt{ 1234567} \\
\texttt{╔═╤═╤═══╗} \\
\texttt{║◌│◌│◌◌◌║ 1} \\
\texttt{║◌├─┤◌◌P║ 2} \\
\texttt{╟─┤◌├─┬─╢ 3} \\
\texttt{║◌│◌│C│◌║ 4} \\
\texttt{╟─┼─┘◌│◌║ 5} \\
\texttt{║◌│L◌┌┴─╢ 6} \\
\texttt{║◌│◌◌│◌◌║ 7} \\
\texttt{╚═╧══╧══╝} \\
\texttt{```} \\
\texttt{Сейчас вы играете раунд 7 из максимум 12.} \\
\texttt{Вы вместе накопили 2 из 8 штрафов.} \\
\texttt{Другой игрок отправил это сообщение:} \\
\\ 
\\ 
\texttt{"Я передвинул L на (3,6). Теперь осталось C на (5,4). Двигаем C. У меня там свободно, проверь у себя."} \\
\\ 
\\ 
\texttt{Какая твоя следующая команда?} \\
            }
        }
    }
    & \\ \\

    \theutterance \stepcounter{utterance}  
    & & & \multicolumn{4}{p{0.6\linewidth}}{
        \cellcolor[rgb]{0.9,0.9,0.9}{
            \makecell[{{p{\linewidth}}}]{
                \texttt{\tiny{[GM$\langle$P2]}}
                \texttt{СКАЗАТЬ: Я уже передвинул C на (5,4). У меня все объекты на местах: P на (7,2), L на (3,6), C на (5,4). Готово?} \\
            }
        }
    }
     \\ \\

    \theutterance \stepcounter{utterance}  
    & & & \multicolumn{2}{p{0.3\linewidth}}{
        \cellcolor[rgb]{0.9,0.9,0.9}{
            \makecell[{{p{\linewidth}}}]{
                \texttt{\tiny{[GM$|$GM]}}
                \texttt{СКАЗАТЬ: Я уже передвинул C на (5,4). У меня все объекты на местах: P на (7,2), L на (3,6), C на (5,4). Готово?} \\
            }
        }
    }
    & & \\ \\

    \theutterance \stepcounter{utterance}  
    & & \multicolumn{4}{p{0.6\linewidth}}{
        \cellcolor[rgb]{0.9,0.9,0.9}{
            \makecell[{{p{\linewidth}}}]{
                \texttt{\tiny{[P1$\langle$GM]}}
                \texttt{Ваше сообщение было передано другому игроку.} \\
\\ 
\\ 
\texttt{Сейчас вы играете раунд 7 из максимум 12.} \\
\texttt{Вы вместе накопили 2 из 8 штрафов.} \\
\texttt{Другой игрок отправил это сообщение:} \\
\\ 
\\ 
\texttt{"Я уже передвинул C на (5,4). У меня все объекты на местах: P на (7,2), L на (3,6), C на (5,4). Готово?"} \\
\\ 
\\ 
\texttt{Какая твоя следующая команда?} \\
            }
        }
    }
    & \\ \\

    \theutterance \stepcounter{utterance}  
    & \multicolumn{4}{p{0.6\linewidth}}{
        \cellcolor[rgb]{0.9,0.9,0.9}{
            \makecell[{{p{\linewidth}}}]{
                \texttt{\tiny{[P1$\rangle$GM]}}
                \texttt{ДВИГАТЬСЯ: C, (5, 4)} \\
            }
        }
    }
    & & \\ \\

    \theutterance \stepcounter{utterance}  
    & & & \multicolumn{2}{p{0.3\linewidth}}{
        \cellcolor[rgb]{0.9,0.9,0.9}{
            \makecell[{{p{\linewidth}}}]{
                \texttt{\tiny{[GM$|$GM]}}
                \texttt{ДВИГАТЬСЯ: C, (5, 4)} \\
            }
        }
    }
    & & \\ \\

    \theutterance \stepcounter{utterance}  
    & & & \multicolumn{2}{p{0.3\linewidth}}{
        \cellcolor[rgb]{0.9,0.9,0.9}{
            \makecell[{{p{\linewidth}}}]{
                \texttt{\tiny{[GM$|$GM]}}
                \texttt{C moved to (5, 4)} \\
            }
        }
    }
    & & \\ \\

    \theutterance \stepcounter{utterance}  
    & & \multicolumn{4}{p{0.6\linewidth}}{
        \cellcolor[rgb]{0.9,0.9,0.9}{
            \makecell[{{p{\linewidth}}}]{
                \texttt{\tiny{[GM$\rangle$P2]}}
                \texttt{Ваше сообщение было передано другому игроку.} \\
\\ 
\\ 
\texttt{Сейчас вы играете раунд 8 из максимум 12.} \\
\texttt{Вы вместе накопили 2 из 8 штрафов.} \\
\texttt{Другой игрок переместил объект на своей сетке.} \\
\\ 
\\ 
\texttt{Какая твоя следующая команда?} \\
            }
        }
    }
    & \\ \\

    \theutterance \stepcounter{utterance}  
    & & & \multicolumn{4}{p{0.6\linewidth}}{
        \cellcolor[rgb]{0.9,0.9,0.9}{
            \makecell[{{p{\linewidth}}}]{
                \texttt{\tiny{[GM$\langle$P2]}}
                \texttt{СКАЗАТЬ: Жду твоего подтверждения. У меня все три объекта на целевых позициях. Готово?} \\
            }
        }
    }
     \\ \\

    \theutterance \stepcounter{utterance}  
    & & & \multicolumn{2}{p{0.3\linewidth}}{
        \cellcolor[rgb]{0.9,0.9,0.9}{
            \makecell[{{p{\linewidth}}}]{
                \texttt{\tiny{[GM$|$GM]}}
                \texttt{СКАЗАТЬ: Жду твоего подтверждения. У меня все три объекта на целевых позициях. Готово?} \\
            }
        }
    }
    & & \\ \\

    \theutterance \stepcounter{utterance}  
    & & \multicolumn{4}{p{0.6\linewidth}}{
        \cellcolor[rgb]{0.9,0.9,0.9}{
            \makecell[{{p{\linewidth}}}]{
                \texttt{\tiny{[P1$\langle$GM]}}
                \texttt{Успешно перемещён объект "C" в (5,4). Ваша обновлённая сетка выглядит так:} \\
\texttt{```} \\
\texttt{ 1234567} \\
\texttt{╔═╤═╤═══╗} \\
\texttt{║◌│◌│◌◌◌║ 1} \\
\texttt{║◌├─┤◌◌P║ 2} \\
\texttt{╟─┤◌├─┬─╢ 3} \\
\texttt{║◌│◌│C│◌║ 4} \\
\texttt{╟─┼─┘◌│◌║ 5} \\
\texttt{║◌│L◌┌┴─╢ 6} \\
\texttt{║◌│◌◌│◌◌║ 7} \\
\texttt{╚═╧══╧══╝} \\
\texttt{```} \\
\texttt{Сейчас вы играете раунд 8 из максимум 12.} \\
\texttt{Вы вместе накопили 2 из 8 штрафов.} \\
\texttt{Другой игрок отправил это сообщение:} \\
\\ 
\\ 
\texttt{"Жду твоего подтверждения. У меня все три объекта на целевых позициях. Готово?"} \\
\\ 
\\ 
\texttt{Какая твоя следующая команда?} \\
            }
        }
    }
    & \\ \\

    \theutterance \stepcounter{utterance}  
    & \multicolumn{4}{p{0.6\linewidth}}{
        \cellcolor[rgb]{0.9,0.9,0.9}{
            \makecell[{{p{\linewidth}}}]{
                \texttt{\tiny{[P1$\rangle$GM]}}
                \texttt{СКАЗАТЬ: готово!} \\
            }
        }
    }
    & & \\ \\

    \theutterance \stepcounter{utterance}  
    & & & \multicolumn{2}{p{0.3\linewidth}}{
        \cellcolor[rgb]{0.9,0.9,0.9}{
            \makecell[{{p{\linewidth}}}]{
                \texttt{\tiny{[GM$|$GM]}}
                \texttt{СКАЗАТЬ: готово!} \\
            }
        }
    }
    & & \\ \\

    \theutterance \stepcounter{utterance}  
    & & \multicolumn{4}{p{0.6\linewidth}}{
        \cellcolor[rgb]{0.9,0.9,0.9}{
            \makecell[{{p{\linewidth}}}]{
                \texttt{\tiny{[GM$\rangle$P2]}}
                \texttt{Ваше сообщение было передано другому игроку.} \\
\\ 
\\ 
\texttt{Сейчас вы играете раунд 9 из максимум 12.} \\
\texttt{Вы вместе накопили 2 из 8 штрафов.} \\
\texttt{Другой игрок отправил это сообщение:} \\
\\ 
\\ 
\texttt{"готово!"} \\
\\ 
\\ 
\texttt{Какая твоя следующая команда?} \\
            }
        }
    }
    & \\ \\

    \theutterance \stepcounter{utterance}  
    & & & \multicolumn{4}{p{0.6\linewidth}}{
        \cellcolor[rgb]{0.9,0.9,0.9}{
            \makecell[{{p{\linewidth}}}]{
                \texttt{\tiny{[GM$\langle$P2]}}
                \texttt{СКАЗАТЬ: готово!} \\
            }
        }
    }
     \\ \\

    \theutterance \stepcounter{utterance}  
    & & & \multicolumn{2}{p{0.3\linewidth}}{
        \cellcolor[rgb]{0.9,0.9,0.9}{
            \makecell[{{p{\linewidth}}}]{
                \texttt{\tiny{[GM$|$GM]}}
                \texttt{СКАЗАТЬ: готово!} \\
            }
        }
    }
    & & \\ \\

    \theutterance \stepcounter{utterance}  
    & & & \multicolumn{2}{p{0.3\linewidth}}{
        \cellcolor[rgb]{0.9,0.9,0.9}{
            \makecell[{{p{\linewidth}}}]{
                \texttt{\tiny{[GM$|$GM]}}
                \texttt{true} \\
            }
        }
    }
    & & \\ \\

    \theutterance \stepcounter{utterance}  
    & & & \multicolumn{2}{p{0.3\linewidth}}{
        \cellcolor[rgb]{0.9,0.9,0.9}{
            \makecell[{{p{\linewidth}}}]{
                \texttt{\tiny{[GM$|$GM]}}
                \texttt{Initial states:} \\
\texttt{```} \\
\texttt{Player 1:} \\
\texttt{ 1234567} \\
\texttt{╔═╤═╤═══╗} \\
\texttt{║L│◌│◌◌◌║ 1} \\
\texttt{║C├─┤◌◌◌║ 2} \\
\texttt{╟─┤◌├─┬─╢ 3} \\
\texttt{║◌│◌│◌│◌║ 4} \\
\texttt{╟─┼─┘◌│◌║ 5} \\
\texttt{║◌│◌◌┌┴─╢ 6} \\
\texttt{║◌│◌◌│◌P║ 7} \\
\texttt{╚═╧══╧══╝} \\
\\ 
\\ 
\texttt{Player 2:} \\
\texttt{ 1234567} \\
\texttt{╔═╤═╤═══╗} \\
\texttt{║◌│◌│◌◌◌║ 1} \\
\texttt{║◌├─┤◌L◌║ 2} \\
\texttt{╟─┤P├─┬─╢ 3} \\
\texttt{║◌│◌│◌│◌║ 4} \\
\texttt{╟─┼─┘C│◌║ 5} \\
\texttt{║◌│◌◌┌┴─╢ 6} \\
\texttt{║◌│◌◌│◌◌║ 7} \\
\texttt{╚═╧══╧══╝} \\
\\ 
\texttt{```} \\
            }
        }
    }
    & & \\ \\

    \theutterance \stepcounter{utterance}  
    & & & \multicolumn{2}{p{0.3\linewidth}}{
        \cellcolor[rgb]{0.9,0.9,0.9}{
            \makecell[{{p{\linewidth}}}]{
                \texttt{\tiny{[GM$|$GM]}}
                \texttt{End states:} \\
\texttt{```} \\
\texttt{Player 1:} \\
\texttt{ 1234567} \\
\texttt{╔═╤═╤═══╗} \\
\texttt{║◌│◌│◌◌◌║ 1} \\
\texttt{║◌├─┤◌◌P║ 2} \\
\texttt{╟─┤◌├─┬─╢ 3} \\
\texttt{║◌│◌│C│◌║ 4} \\
\texttt{╟─┼─┘◌│◌║ 5} \\
\texttt{║◌│L◌┌┴─╢ 6} \\
\texttt{║◌│◌◌│◌◌║ 7} \\
\texttt{╚═╧══╧══╝} \\
\\ 
\\ 
\texttt{Player 2:} \\
\texttt{ 1234567} \\
\texttt{╔═╤═╤═══╗} \\
\texttt{║◌│◌│◌◌◌║ 1} \\
\texttt{║◌├─┤◌◌P║ 2} \\
\texttt{╟─┤◌├─┬─╢ 3} \\
\texttt{║◌│◌│C│◌║ 4} \\
\texttt{╟─┼─┘◌│◌║ 5} \\
\texttt{║◌│L◌┌┴─╢ 6} \\
\texttt{║◌│◌◌│◌◌║ 7} \\
\texttt{╚═╧══╧══╝} \\
\\ 
\texttt{```} \\
            }
        }
    }
    & & \\ \\

    \theutterance \stepcounter{utterance}  
    & & & \multicolumn{2}{p{0.3\linewidth}}{
        \cellcolor[rgb]{0.9,0.9,0.9}{
            \makecell[{{p{\linewidth}}}]{
                \texttt{\tiny{[GM$|$GM]}}
                \texttt{* success: True} \\
\texttt{* lose: False} \\
\texttt{* aborted: False} \\
\texttt{{-}{-}{-}{-}{-}{-}{-}} \\
\texttt{* Shifts: 2.00} \\
\texttt{* Max Shifts: 4.00} \\
\texttt{* Min Shifts: 2.00} \\
\texttt{* End Distance Sum: 0.00} \\
\texttt{* Init Distance Sum: 15.76} \\
\texttt{* Expected Distance Sum: 12.57} \\
\texttt{* Penalties: 2.00} \\
\texttt{* Max Penalties: 8.00} \\
\texttt{* Rounds: 9.00} \\
\texttt{* Max Rounds: 12.00} \\
\texttt{* Object Count: 3.00} \\
            }
        }
    }
    & & \\ \\

\end{supertabular}
}

\end{document}
